% =============================================================================
% Sprint 2 — Intégration continue et approvisionnement de l'infrastructure
% TicketManager — aligné sur le dépôt réel (terraform/, azure-pipelines.yml)
% Captures d'écran : placer les PNG dans le dossier img/ du rapport LaTeX
% =============================================================================

\section[Sprint 2 : Implémentation de l'intégration continue et de l'infrastructure]%
{Sprint 2 : Implémentation de l'intégration continue et de l'infrastructure}

Ce sprint vise à industrialiser le cycle de vie de la plateforme \textit{TicketManager} en
automatisant la validation du code, la construction des artefacts et l'approvisionnement de
l'infrastructure sur \textbf{Azure Stack Hub}. Il s'appuie sur une approche \textit{Infrastructure
as Code} pour le réseau et les machines virtuelles, complétée par un pipeline d'intégration
continue défini de manière déclarative.


% ---------------------------------------------------------------------------
\subsection{Objectif du Sprint 2}
% ---------------------------------------------------------------------------

Le Sprint~2 prolonge le socle applicatif réalisé au Sprint~1. Réalisé en \textbf{autonomie}
(développement, intégration et infrastructure), il met l'accent sur
\textbf{l'industrialisation} et \textbf{l'approvisionnement cloud}. L'objectif est de rendre
reproductibles la validation du code, la production d'artefacts conteneurisés et le déploiement
de l'infrastructure cible sur Azure Stack Hub.

\begin{table}[H]
\centering
\footnotesize
\renewcommand{\arraystretch}{1.1}
\setlength{\tabcolsep}{6pt}
\begin{tabular}{|p{2.8cm}|p{9.2cm}|}
\hline
\textbf{Axe} & \textbf{Objectif} \\
\hline
Infrastructure & Définir et provisionner le réseau, la sécurité et les machines virtuelles
                 par Terraform (\textit{Infrastructure as Code}) sur Azure Stack Hub \\
\hline
Intégration continue & Mettre en place un pipeline Azure DevOps déclaratif (YAML) déclenché
                       à chaque évolution du code source \\
\hline
Qualité et sécurité & Automatiser les tests, l'analyse statique (SonarQube), l'analyse des
                      dépendances (Snyk) et le scan des images (Trivy) \\
\hline
Conteneurisation & Construire des images Docker multi-étapes pour le backend et le frontend,
                   puis les publier dans un registre privé (Nexus) \\
\hline
Préparation au déploiement & Préparer le sous-réseau Kubernetes et la chaîne d'outils DevOps
                             (registre, analyse) en vue du déploiement applicatif au sprint suivant \\
\hline
\end{tabular}
\caption{Objectifs du Sprint~2 par axe}
\label{tab:s2-objectifs}
\end{table}

\paragraph{Résultat attendu.}
À l'issue de ce sprint, la solution \textit{TicketManager} dispose des éléments suivants :
\begin{itemize}
  \item une infrastructure réseau documentée et reproductible (VNet, sous-réseaux, NSG,
        machines de gestion et de base de données) ;
  \item un pipeline CI/CD exécutable sur agent auto-hébergé, couvrant compilation, tests,
        analyse qualité, construction d'images, scan de vulnérabilités et publication vers Nexus ;
  \item une étape Terraform intégrée au pipeline, avec contrôle manuel avant application sur
        la branche principale ;
  \item un socle prêt pour le déploiement du cluster Kubernetes et des services applicatifs
        au sprint ultérieur.
\end{itemize}

\paragraph{Périmètre.}
Le tableau~\ref{tab:s2-perimetre} précise le contenu du Sprint~2 par rapport aux livraisons
précédentes et futures.

\begin{table}[H]
\centering
\footnotesize
\renewcommand{\arraystretch}{1.1}
\setlength{\tabcolsep}{6pt}
\begin{tabular}{|p{5.5cm}|c|c|}
\hline
\textbf{Élément} & \textbf{Sprint~2} & \textbf{Sprint ultérieur} \\
\hline
Socle applicatif (backend, frontend, MongoDB locale) & — & Sprint~1 \\
\hline
Terraform — réseau, NSG, VMs (gestion, base de données) & Oui & — \\
\hline
Pipeline Azure DevOps (build, tests, qualité, Docker) & Oui & — \\
\hline
Registre Docker privé (Nexus) & Oui & — \\
\hline
Déploiement cluster Kubernetes (AKS Engine) & Préparation & Sprint~3 \\
\hline
Déploiement applicatif en pods (frontend, backend) & — & Sprint~3 \\
\hline
Déploiement continu (CD) vers production & — & Sprint~3 \\
\hline
\end{tabular}
\caption{Périmètre fonctionnel et infrastructurel du Sprint~2}
\label{tab:s2-perimetre}
\end{table}


% ---------------------------------------------------------------------------
\subsection{Backlog du Sprint 2}
% ---------------------------------------------------------------------------

Le backlog du Sprint~2 est organisé en trois epics : infrastructure cloud, chaîne
d'intégration continue et outillage DevOps. L'acteur principal des user stories est le
\textbf{développeur}, seul responsable de la conception, de l'implémentation et de la mise en
œuvre DevOps du projet. Le tableau~\ref{tab:s2-backlog-consolide} en présente une vue
consolidée ; les tableaux~\ref{tab:s2-backlog-infra} et~\ref{tab:s2-backlog-cicd} détaillent
les livrables par epic.

\subsubsection{Epic — Infrastructure cloud (Terraform)}

\begin{table}[H]
\centering
\footnotesize
\renewcommand{\arraystretch}{0.9}
\setlength{\tabcolsep}{4pt}
\begin{tabular}{|c|p{5.2cm}|c|p{4.2cm}|c|}
\hline
\textbf{ID} & \textbf{User story} & \textbf{Priorité} & \textbf{Livrable} & \textbf{Statut} \\
\hline
S2-I01 & En tant que développeur, je provisionne l'infrastructure par Terraform sur Azure Stack Hub & Must & Cycle init, validate, plan, apply & $\checkmark$ \\
\hline
S2-I02 & En tant que développeur, je segmente le réseau en sous-réseaux dédiés (administration, base de données, Kubernetes) & Must & VNet \texttt{10.1.0.0/16}, 3 subnets & $\checkmark$ \\
\hline
S2-I03 & En tant que développeur, je configure les groupes de sécurité réseau par sous-réseau & Must & NSG par sous-réseau & $\checkmark$ \\
\hline
S2-I04 & En tant que développeur, je dispose d'une VM d'administration pour orchestrer les déploiements & Must & VM-gestion (IP publique + privée) & $\checkmark$ \\
\hline
S2-I05 & En tant que développeur, je déploie une VM hébergeant MongoDB en environnement cloud & Must & VM-base de données (MongoDB) & $\checkmark$ \\
\hline
S2-I06 & En tant que développeur, j'associe chaque VM à une interface réseau sécurisée & Must & NIC + association NSG & $\checkmark$ \\
\hline
S2-I07 & En tant que développeur, je prépare le sous-réseau pour le futur cluster Kubernetes & Should & Subnet AKS, règles NSG associées & $\checkmark$ \\
\hline
\end{tabular}
\caption{Backlog Sprint~2 — Epic infrastructure cloud}
\label{tab:s2-backlog-infra}
\end{table}

\subsubsection{Epic — Intégration continue (Azure DevOps)}

\begin{table}[H]
\centering
\footnotesize
\renewcommand{\arraystretch}{0.9}
\setlength{\tabcolsep}{4pt}
\begin{tabular}{|c|p{5.2cm}|c|p{4.2cm}|c|}
\hline
\textbf{ID} & \textbf{User story} & \textbf{Priorité} & \textbf{Livrable} & \textbf{Statut} \\
\hline
S2-C01 & En tant que développeur, je centralise le code source sur Azure Repos & Must & Dépôt distant, synchronisation Git & $\checkmark$ \\
\hline
S2-C02 & En tant que développeur, j'exécute le pipeline sur un agent auto-hébergé & Must & Pool auto-hébergé \textit{ticketmanager-pool} & $\checkmark$ \\
\hline
S2-C03 & En tant que développeur, je déclenche automatiquement le pipeline à chaque push & Must & Fichier \texttt{azure-pipelines.yml}, branches \texttt{main}/\texttt{develop} & $\checkmark$ \\
\hline
S2-C04 & En tant que développeur, je compile backend et frontend à chaque livraison & Must & Stage Build (NestJS, Next.js) & $\checkmark$ \\
\hline
S2-C05 & En tant que développeur, j'exécute les tests automatisés et publie la couverture & Must & Stage Test (unitaires, intégration, couverture) & $\checkmark$ \\
\hline
S2-C06 & En tant que développeur, j'intègre l'analyse SonarQube dans le pipeline & Must & Stage SonarQube & $\checkmark$ \\
\hline
S2-C07 & En tant que développeur, j'analyse les vulnérabilités des dépendances avec Snyk & Should & Stage Snyk & $\checkmark$ \\
\hline
S2-C08 & En tant que développeur, je construis des images Docker prêtes pour la production & Must & Dockerfiles multi-étapes, stage Docker & $\checkmark$ \\
\hline
S2-C09 & En tant que développeur, je scanne les images Docker avant publication & Must & Stage Trivy & $\checkmark$ \\
\hline
S2-C10 & En tant que développeur, je publie les images validées dans un registre privé & Must & Stage Push vers Nexus & $\checkmark$ \\
\hline
S2-C11 & En tant que développeur, j'automatise Terraform dans le pipeline avec validation manuelle & Must & Stage Terraform (approbation sur \texttt{main}) & $\checkmark$ \\
\hline
\end{tabular}
\caption{Backlog Sprint~2 — Epic intégration continue}
\label{tab:s2-backlog-cicd}
\end{table}

\subsubsection{Vue consolidée du backlog}

\begin{table}[H]
\centering
\footnotesize
\renewcommand{\arraystretch}{0.85}
\setlength{\tabcolsep}{3pt}
\begin{tabular}{|c|p{6.8cm}|c|p{2.8cm}|c|}
\hline
\textbf{ID} & \textbf{User story} & \textbf{Priorité} & \textbf{Epic} & \textbf{Statut} \\
\hline
S2-01 & En tant que développeur, je provisionne l'infrastructure réseau et compute par Terraform & Must & Infrastructure & $\checkmark$ \\
\hline
S2-02 & En tant que développeur, je déploie les VMs de gestion et de base de données & Must & Infrastructure & $\checkmark$ \\
\hline
S2-03 & En tant que développeur, je configure les NSG et les interfaces réseau & Must & Infrastructure & $\checkmark$ \\
\hline
S2-04 & En tant que développeur, je connecte mon dépôt local à Azure Repos & Must & CI/CD & $\checkmark$ \\
\hline
S2-05 & En tant que développeur, je configure un agent auto-hébergé pour le pipeline & Must & CI/CD & $\checkmark$ \\
\hline
S2-06 & En tant que développeur, j'automatise la compilation backend et frontend & Must & CI/CD & $\checkmark$ \\
\hline
S2-07 & En tant que développeur, j'exécute les tests unitaires et d'intégration à chaque build & Must & CI/CD & $\checkmark$ \\
\hline
S2-08 & En tant que développeur, j'intègre SonarQube dans le pipeline & Must & Outillage & $\checkmark$ \\
\hline
S2-09 & En tant que développeur, j'analyse les vulnérabilités des dépendances avec Snyk & Should & Outillage & $\checkmark$ \\
\hline
S2-10 & En tant que développeur, je construis les images Docker multi-étapes & Must & CI/CD & $\checkmark$ \\
\hline
S2-11 & En tant que développeur, je scanne les images Docker avec Trivy & Must & Outillage & $\checkmark$ \\
\hline
S2-12 & En tant que développeur, je publie les images vers Nexus & Must & Outillage & $\checkmark$ \\
\hline
S2-13 & En tant que développeur, j'intègre Terraform au pipeline avec contrôle manuel & Must & CI/CD & $\checkmark$ \\
\hline
S2-14 & En tant que développeur, je prépare le sous-réseau et la configuration AKS Engine & Should & Infrastructure & En cours \\
\hline
\end{tabular}
\caption{Backlog consolidé du Sprint~2}
\label{tab:s2-backlog-consolide}
\end{table}

\paragraph{Hors périmètre Sprint~2.}
Les éléments suivants sont reportés au sprint suivant :
\begin{itemize}
  \item déploiement effectif du cluster Kubernetes via AKS Engine ;
  \item déploiement des pods applicatifs (frontend, backend) en production ;
  \item mise en place du déploiement continu (CD) automatisé vers le cluster ;
  \item fonctionnalités métier additionnelles (réinitialisation mot de passe, profil utilisateur).
\end{itemize}


% ---------------------------------------------------------------------------
\subsection{Approvisionnement de l'infrastructure cloud}
% ---------------------------------------------------------------------------

\subsubsection{Infrastructure as Code avec Terraform}

Terraform est un outil d'Infrastructure as Code (IaC) édité par HashiCorp. Il permet de
décrire de façon déclarative l'état cible des ressources cloud, sans détailler manuellement
chaque opération de création. Dans le cadre de TicketManager, Terraform est utilisé avec le
fournisseur \textit{azurestack}, adapté au environnement Azure Stack Hub de l'organisme
d'hébergement.

Le cycle de déploiement repose sur quatre commandes complémentaires :

\begin{itemize}
  \item \textbf{Initialisation} (\texttt{terraform init}) : prépare l'environnement de travail,
        télécharge le fournisseur et prépare le backend local d'état ;
  \item \textbf{Validation} (\texttt{terraform validate}) : contrôle la cohérence syntaxique
        des fichiers de configuration avant toute exécution ;
  \item \textbf{Planification} (\texttt{terraform plan}) : produit un plan d'exécution
        détaillant les ressources à créer, modifier ou supprimer ;
  \item \textbf{Application} (\texttt{terraform apply}) : matérialise le plan en respectant
        l'ordre des dépendances entre ressources.
\end{itemize}

La figure~\ref{fig:tf-init} illustre la réussite de l'étape d'initialisation.
La figure~\ref{fig:tf-validate} confirme la validité de la configuration.
La figure~\ref{fig:tf-plan} présente le plan d'exécution généré.
La figure~\ref{fig:tf-apply} montre le résultat de l'application, incluant les adresses
réseau attribuées aux machines virtuelles.

\begin{figure}[htbp]
    \centering
    \includegraphics[width=0.85\textwidth]{img/terraform-init.png}
    \caption{Résultat de l'initialisation Terraform}
    \label{fig:tf-init}
\end{figure}

\begin{figure}[htbp]
    \centering
    \includegraphics[width=0.85\textwidth]{img/terraform-validate.png}
    \caption{Validation réussie de la configuration Terraform}
    \label{fig:tf-validate}
\end{figure}

\begin{figure}[htbp]
    \centering
    \includegraphics[width=0.85\textwidth]{img/terraform-plan.png}
    \caption{Plan d'exécution Terraform}
    \label{fig:tf-plan}
\end{figure}

\begin{figure}[htbp]
    \centering
    \includegraphics[width=0.85\textwidth]{img/terraform-apply.png}
    \caption{Résultat de l'application Terraform (adresses IP et sorties)}
    \label{fig:tf-apply}
\end{figure}

L'utilisation de Terraform apporte plusieurs bénéfices pour un projet académique de PFE :
reproductibilité des déploiements, traçabilité via le fichier d'état (\texttt{terraform.tfstate}),
réduction des erreurs manuelles et homogénéité entre les environnements de test et de
démonstration.


\subsubsection{Architecture de l'infrastructure}

L'architecture retenue pour TicketManager repose sur une organisation hybride : une couche
réseau et compute provisionnée par Terraform, complétée par un déploiement Kubernetes prévu
via AKS Engine sur un sous-réseau dédié. La figure~\ref{fig:arch-infra} en donne une vue
d'ensemble.

Toutes les ressources sont regroupées dans un même groupe de ressources Azure Stack Hub,
attribué par l'encadrant du projet. L'infrastructure comprend les composants suivants.

\paragraph{Réseau virtuel et sous-réseaux.}
Un réseau virtuel (\texttt{10.1.0.0/16}) structure la connectivité interne. Trois sous-réseaux
sont définis :
\begin{itemize}
  \item \textbf{Administration} (\texttt{10.1.1.0/24}) : hébergement de la machine de gestion ;
  \item \textbf{Base de données} (\texttt{10.1.2.0/24}) : hébergement de MongoDB, sans
        exposition directe à Internet ;
  \item \textbf{Kubernetes} (\texttt{10.1.3.0/24}) : réservé au futur cluster AKS Engine.
\end{itemize}
La figure~\ref{fig:vnet-subnets} illustre cette segmentation.

\paragraph{Machines virtuelles.}
Deux machines virtuelles Linux sont déployées par Terraform (et non trois) :
\begin{itemize}
  \item \textbf{VM-gestion} : point d'administration pour l'exploitation de l'infrastructure,
        l'exécution d'AKS Engine et les opérations de déploiement ;
  \item \textbf{VM-base de données} : héberge MongoDB pour la persistance des tickets et des
        utilisateurs en environnement cloud.
\end{itemize}
La figure~\ref{fig:vms} présente ces machines dans le portail Azure Stack Hub.

\paragraph{Groupes de sécurité réseau.}
Trois groupes de sécurité réseau (NSG) filtrent le trafic par rôle : administration (accès
SSH), base de données (MongoDB depuis les sous-réseaux autorisés) et cluster Kubernetes
(trafic interne au réseau virtuel). La figure~\ref{fig:nsg} les représente.

\paragraph{Interfaces réseau.}
Chaque machine virtuelle est rattachée à une interface réseau propre, intégrée au réseau
virtuel et associée au NSG correspondant. La figure~\ref{fig:nics} montre ces associations.

\paragraph{Cluster Kubernetes (perspective).}
Le cluster Kubernetes n'est pas créé par Terraform lui-même : le module réseau prépare
uniquement le sous-réseau et les règles associées. Le déploiement du plan de contrôle et des
nœuds workers est confié à AKS Engine, outil adapté à Azure Stack Hub. Les services applicatifs
conteneurisés (frontend Next.js, backend NestJS) ainsi que les outils DevOps (registre Nexus,
analyse SonarQube, supervision) seront déployés sous forme de pods dans ce cluster.

\begin{figure}[htbp]
    \centering
    \includegraphics[width=0.9\textwidth]{img/architecture-infra-ticketmanager.png}
    \caption{Architecture globale de l'infrastructure TicketManager}
    \label{fig:arch-infra}
\end{figure}

\begin{figure}[htbp]
    \centering
    \includegraphics[width=0.85\textwidth]{img/vnet-subnets.png}
    \caption{Réseau virtuel et sous-réseaux}
    \label{fig:vnet-subnets}
\end{figure}

\begin{figure}[htbp]
    \centering
    \includegraphics[width=0.85\textwidth]{img/vms-azurestack.png}
    \caption{Machines virtuelles déployées sur Azure Stack Hub}
    \label{fig:vms}
\end{figure}

\begin{figure}[htbp]
    \centering
    \includegraphics[width=0.85\textwidth]{img/nsg-azurestack.png}
    \caption{Groupes de sécurité réseau}
    \label{fig:nsg}
\end{figure}

\begin{figure}[htbp]
    \centering
    \includegraphics[width=0.85\textwidth]{img/nics-azurestack.png}
    \caption{Interfaces réseau des machines virtuelles}
    \label{fig:nics}
\end{figure}


% ---------------------------------------------------------------------------
\subsection{Présentation d'Azure DevOps}
% ---------------------------------------------------------------------------

Azure DevOps est la plateforme collaborative retenue pour centraliser le code source, planifier
le travail et automatiser l'intégration continue. La figure~\ref{fig:azure-devops-ui} présente
son interface principale.

Les modules mobilisés dans TicketManager sont les suivants :
\begin{itemize}
  \item \textbf{Azure Boards} : planification agile, suivi des sprints et des user stories ;
  \item \textbf{Azure Repos} : gestion de versions Git, branches et revues de code ;
  \item \textbf{Azure Pipelines} : exécution automatisée du pipeline CI/CD défini en YAML.
\end{itemize}


% ---------------------------------------------------------------------------
\subsection{Conception et infrastructure du pipeline CI}
% ---------------------------------------------------------------------------

\subsubsection{Approche déclarative en YAML}

Le pipeline est entièrement décrit dans le fichier \texttt{azure-pipelines.yml}, versionné avec
le code source. Cette approche garantit la traçabilité des modifications, la reproductibilité
des exécutions et la cohérence entre les environnements.

\subsubsection{Pool d'agents auto-hébergés}

Les tâches du pipeline s'exécutent sur un \textbf{agent auto-hébergé}, regroupé dans le pool
\textit{ticketmanager-pool}. Ce choix permet de maîtriser l'environnement d'exécution
(Docker, Node.js, outils d'analyse) et de rester compatible avec l'infrastructure locale ou
Stack Hub. La figure~\ref{fig:agent-pool} illustre la configuration du pool et de l'agent.

La figure~\ref{fig:agent-config} montre la procédure de liaison de l'agent au projet Azure
DevOps.

\subsubsection{Connexion au dépôt distant}

Le code source est synchronisé avec Azure Repos selon le flux Git classique : initialisation
du dépôt local, indexation des modifications, création de commits descriptifs et envoi vers
le dépôt distant. La figure~\ref{fig:azure-repos} présente l'arborescence du projet, incluant
le fichier de pipeline et les répertoires \texttt{backend}, \texttt{frontend} et
\texttt{terraform}.

\begin{figure}[htbp]
    \centering
    \includegraphics[width=0.85\textwidth]{img/azure-devops-interface.png}
    \caption{Interface principale d'Azure DevOps}
    \label{fig:azure-devops-ui}
\end{figure}

\begin{figure}[htbp]
    \centering
    \includegraphics[width=0.85\textwidth]{img/agent-pool.png}
    \caption{Configuration du pool d'agents auto-hébergés}
    \label{fig:agent-pool}
\end{figure}

\begin{figure}[htbp]
    \centering
    \includegraphics[width=0.85\textwidth]{img/agent-configuration.png}
    \caption{Liaison de l'agent auto-hébergé au projet}
    \label{fig:agent-config}
\end{figure}

\begin{figure}[htbp]
    \centering
    \includegraphics[width=0.85\textwidth]{img/azure-repos.png}
    \caption{Dépôt Azure Repos — structure du projet TicketManager}
    \label{fig:azure-repos}
\end{figure}


% ---------------------------------------------------------------------------
\subsection{Étapes du pipeline d'intégration continue}
% ---------------------------------------------------------------------------

Le pipeline est organisé en \textbf{stages} successifs, déclenchés automatiquement à chaque
push sur les branches \texttt{main} et \texttt{develop}. La figure~\ref{fig:pipeline-overview}
en présente la vue d'ensemble.

\begin{figure}[htbp]
    \centering
    \includegraphics[width=0.95\textwidth]{img/pipeline-stages.png}
    \caption{Vue d'ensemble des étapes du pipeline CI/CD}
    \label{fig:pipeline-overview}
\end{figure}


\subsubsection{Compilation des composants applicatifs}

La première étape compile le backend (API NestJS) et le frontend (application Next.js) après
installation des dépendances. Un mécanisme de cache accélère les exécutions suivantes en
réutilisant les modules déjà téléchargés.


\subsubsection{Vérification du code et exécution des tests}

Les tests automatisés constituent un garde-fou essentiel de l'intégration continue.

\paragraph{Tests backend.}
Le backend est couvert par des tests unitaires et des tests d'intégration, exécutés avec
l'outil Jest. Les résultats sont publiés dans Azure DevOps sous forme de rapports JUnit, et
un rapport de couverture de code est généré au format Cobertura. La figure~\ref{fig:tests-backend}
illustre l'exécution réussie de cette étape.

\paragraph{Tests frontend.}
La validation du frontend repose principalement sur la compilation de production et les
contrôles qualité intégrés au pipeline ; des tests unitaires frontend pourront être renforcés
dans une itération ultérieure.

\begin{figure}[htbp]
    \centering
    \includegraphics[width=0.85\textwidth]{img/tests-backend.png}
    \caption{Exécution des tests automatisés du backend}
    \label{fig:tests-backend}
\end{figure}


\subsubsection{Analyse de la qualité du code}

SonarQube est intégré au pipeline pour l'analyse statique du code source backend et frontend.
Il mesure la dette technique, détecte les duplications et évalue la couverture des tests.
La connexion au serveur SonarQube s'effectue via un jeton d'authentification configuré dans
Azure DevOps. La figure~\ref{fig:sonar-connection} présente la liaison entre les deux outils,
et la figure~\ref{fig:sonar-results} un extrait du tableau de bord d'analyse.


\subsubsection{Analyse de sécurité des dépendances}

En parallèle, l'outil Snyk analyse les dépendances npm du backend et du frontend afin
d'identifier les vulnérabilités connues avant la construction des images conteneurisées.


\subsubsection{Construction des images Docker}

Une fois le code validé, deux images Docker sont produites :
\begin{itemize}
  \item \textbf{Image backend} : API NestJS exposée sur le port 3001 ;
  \item \textbf{Image frontend} : application Next.js exposée sur le port 3000.
\end{itemize}

Chaque image est construite selon une approche \textbf{multi-étapes} : une phase de compilation
séparée de l'image d'exécution finale, allégée et adaptée à la production. Les
figures~\ref{fig:dockerfile-frontend} et~\ref{fig:dockerfile-backend} présentent les
Dockerfiles correspondants ; la figure~\ref{fig:docker-build} montre le succès de la
construction dans le pipeline.

\begin{figure}[htbp]
    \centering
    \includegraphics[width=0.85\textwidth]{img/dockerfile-frontend.png}
    \caption{Dockerfile du frontend Next.js}
    \label{fig:dockerfile-frontend}
\end{figure}

\begin{figure}[htbp]
    \centering
    \includegraphics[width=0.85\textwidth]{img/dockerfile-backend.png}
    \caption{Dockerfile du backend NestJS}
    \label{fig:dockerfile-backend}
\end{figure}

\begin{figure}[htbp]
    \centering
    \includegraphics[width=0.85\textwidth]{img/docker-build-pipeline.png}
    \caption{Construction des images Docker dans le pipeline}
    \label{fig:docker-build}
\end{figure}


\subsubsection{Scan de vulnérabilités des images}

L'outil Trivy analyse les images Docker construites afin de détecter les failles de sécurité
de niveau élevé ou critique. Les figures~\ref{fig:trivy-backend} et~\ref{fig:trivy-frontend}
présentent les rapports obtenus pour le backend et le frontend. Les anomalies identifiées
peuvent être corrigées par mise à jour des dépendances avant publication.

\begin{figure}[htbp]
    \centering
    \includegraphics[width=0.85\textwidth]{img/trivy-backend.png}
    \caption{Analyse Trivy — image backend}
    \label{fig:trivy-backend}
\end{figure}

\begin{figure}[htbp]
    \centering
    \includegraphics[width=0.85\textwidth]{img/trivy-frontend.png}
    \caption{Analyse Trivy — image frontend}
    \label{fig:trivy-frontend}
\end{figure}


\subsubsection{Publication des images vers Nexus Repository}

Nexus Repository Manager joue le rôle de \textbf{registre Docker privé}. Il centralise les
images validées par le pipeline et assure leur traçabilité avant déploiement dans le cluster
Kubernetes.

\paragraph{Positionnement dans l'architecture.}
Dans la cible retenue pour TicketManager, Nexus est hébergé au sein du cluster Kubernetes
sous forme de pod dédié, aux côtés des autres outils DevOps (analyse qualité, supervision).
Le registre Docker y est exposé via un service interne au cluster, accessible depuis l'agent
d'intégration continue. Cette organisation évite la dispersion des artefacts et prépare le
déploiement continu des composants applicatifs, qui consommeront les mêmes images depuis le
registre interne.

\paragraph{Configuration.}
Un dépôt Docker de type \textit{hosted} est créé dans Nexus, associé à un compte de service
disposant des droits de lecture et d'écriture. Les figures~\ref{fig:nexus-repo}
et~\ref{fig:nexus-user} illustrent cette configuration initiale.

\paragraph{Intégration pipeline.}
Après authentification au registre, les images backend et frontend sont étiquetées avec
l'identifiant de build puis poussées vers Nexus. Sur la branche principale, une étiquette
stable \texttt{latest} est également appliquée afin de faciliter les déploiements de
démonstration. La figure~\ref{fig:nexus-images} montre les images stockées après exécution
du pipeline.

\begin{figure}[htbp]
    \centering
    \includegraphics[width=0.85\textwidth]{img/sonar-connection.png}
    \caption{Connexion SonarQube — Azure DevOps}
    \label{fig:sonar-connection}
\end{figure}

\begin{figure}[htbp]
    \centering
    \includegraphics[width=0.85\textwidth]{img/sonar-results.png}
    \caption{Résultats d'analyse SonarQube}
    \label{fig:sonar-results}
\end{figure}

\begin{figure}[htbp]
    \centering
    \includegraphics[width=0.85\textwidth]{img/nexus-docker-repo.png}
    \caption{Création du dépôt Docker dans Nexus}
    \label{fig:nexus-repo}
\end{figure}

\begin{figure}[htbp]
    \centering
    \includegraphics[width=0.85\textwidth]{img/nexus-user.png}
    \caption{Compte de service Nexus pour le registre Docker}
    \label{fig:nexus-user}
\end{figure}

\begin{figure}[htbp]
    \centering
    \includegraphics[width=0.85\textwidth]{img/nexus-images.png}
    \caption{Images Docker publiées dans Nexus après exécution du pipeline}
    \label{fig:nexus-images}
\end{figure}


\subsubsection{Provisionnement Terraform automatisé}

En fin de pipeline, une étape dédiée exécute le cycle Terraform sur Azure Stack Hub
(initialisation, validation, planification). Sur la branche principale, une \textbf{validation
manuelle} est requise avant l'application effective, afin de contrôler les modifications
infrastructurelles. La figure~\ref{fig:terraform-pipeline} illustre cette étape dans Azure
Pipelines.

\begin{figure}[htbp]
    \centering
    \includegraphics[width=0.85\textwidth]{img/terraform-pipeline-stage.png}
    \caption{Étape Terraform intégrée au pipeline CI/CD}
    \label{fig:terraform-pipeline}
\end{figure}


% ---------------------------------------------------------------------------
\subsection{Bilan du Sprint 2}
% ---------------------------------------------------------------------------

À l'issue de ce sprint, TicketManager dispose d'une chaîne d'intégration continue couvrant
la compilation, les tests, l'analyse qualité et sécurité, la conteneurisation, la publication
d'artefacts et le provisionnement réseau sur Azure Stack Hub. Cette industrialisation constitue
le socle nécessaire au déploiement applicatif sur le cluster Kubernetes prévu au sprint
suivant, ainsi qu'à la reproductibilité des livraisons dans un contexte académique de PFE.
